% !TeX root = ../CV note.tex
\section{Image Processing}

\subsection{Pixel (Point) Processing}

\subsubsection{Basic Pixel Operations}
Pixel Processing的基础操作主要有以下几个层面：

\begin{itemize}
    \item \textbf{变暗（darken）：}$I(x,y) - 128$
    \item \textbf{变亮（lighten）：}$I(x,y) + 128$
    \item \textbf{降低对比度（low contrast）：}$I(x,y)/2$
    \item \textbf{反色（invert）：}$255 - I(x,y)$
    \item \textbf{灰度化（gray）：}$\mathbf{dot}\bigl([0.3,\ 0.6,\ 0.1],\ I(x,y)\bigr)$，三个权重依次对应 R、G、B 通道，绿色权重最大
\end{itemize}

\subsubsection{Histogram Equalization}
Pixel Processing中较复杂的是直方图均衡化方法，用来调整图像中RGB通道的强度。用意在于调整比较复杂的颜色密度分布变换为均匀分布。这个方法的数学原理是将概率分布变换为均匀分布，寻找一个单调递增的映射函数T（为了保证是可逆的不丢失信息），将原有分布复合通过这个映射之后得到均匀分布。

具体的技巧是，首先把RGB的灰度分布曲线画出来（对每一个颜色通道，统计每种灰度对应的像素个数，然后归一化计算概率分布），得到概率密度函数pdf，积分得到cdf，将CDF的轴拉成255对应灰度值范围，然后将原始灰度数值映射为新的处理过后的灰度值。

\[
s_k = (L - 1)CDF(r_k)
\]

为什么CDF呢？论证的过程需要依照随机变量变换定理，变换的规则是变换前后在同一个足够小区间上的必须是一致的（概率总量守恒）

根据概率论中的随机变量变换定理，变量变换前后在区间 $[r, r+dr]$ 与 $[s, s+ds]$ 内的概率总量必须守恒：$$p_s(s) \vert{}ds\vert{} = p_r(r) \vert{}dr\vert{}$$

设原始图像的灰度值是一个连续随机变量 $r \in [0, L-1]$，其概率密度函数（PDF）为 $p_r(r)$。我们希望寻找一个单调递增的映射函数 $s = T(r)$，使得变换后的新灰度变量 $s$ 在 $[0, L-1]$ 上服从均匀分布，即：$$p_s(s) = \frac{1}{L-1}$$

由于映射函数 $T(r)$ 是单调递增的（以保证像素灰度相对大小顺序不颠倒，避免图像反色），有 $\frac{ds}{dr} > 0$，因此：$$p_s(s) = p_r(r) \cdot \frac{dr}{ds}$$3. 导出映射函数：累积分布函数（CDF）将要求的均匀分布 $p_s(s) = \frac{1}{L-1}$ 

代入上式：$$\frac{1}{L-1} = p_r(r) \cdot \frac{dr}{ds} \implies \frac{ds}{dr} = (L-1) p_r(r)$$对微分方程两边关于 $r$ 进行积分：$$s = T(r) = (L-1) \int_0^r p_r(w) \, dw = (L-1) \cdot F_r(r)$$其中 $F_r(r)$ 正是原始灰度的累积分布函数（CDF）。

值得注意的一个拓展：随机变量变换定理也可以倒过来用，C语言标准库中，使用标准随机数在y轴上面取值，反函数映射到x轴上，即可得到关于x的均匀随机分布。

\subsection{Filtering}

\subsubsection{Linear Filters}
\paragraph{Classification of Linear Filters}
图像滤波常常用于图像的模糊，锐化，特征提取等算法。线性滤波器就是根据不同的过滤原理实现卷积算子，是一个3*3的矩阵kernel，然后在图像周围补0,扫描整个图像。
\paragraph{Box Filter}
卷积核为全 1 矩阵除以元素个数，即邻域算术平均，得到低通模糊：
$$K_{\text{box}} = \frac{1}{9}\begin{bmatrix} 1 & 1 & 1 \\ 1 & 1 & 1 \\ 1 & 1 & 1 \end{bmatrix}$$

\paragraph{Gaussian Filter}
卷积核由二维高斯函数采样并归一化，中心权重大、四周小，同属低通滤波（Low Pass Filtering）：
$$K_{\text{gauss}} = \frac{1}{16}\begin{bmatrix} 1 & 2 & 1 \\ 2 & 4 & 2 \\ 1 & 2 & 1 \end{bmatrix}$$

\paragraph{Sharpening Filter}
单位冲激核减去均值核，得到高通核：
$$K_{\text{sharp}} = \begin{bmatrix} 0 & 0 & 0 \\ 0 & 2 & 0 \\ 0 & 0 & 0 \end{bmatrix} - \frac{1}{9}\begin{bmatrix} 1 & 1 & 1 \\ 1 & 1 & 1 \\ 1 & 1 & 1 \end{bmatrix}$$
等价于在原图上叠加被放大的高频成分：
$$\text{Sharpening: } I + \lambda\,(I - I_{\text{blur}})$$
其中 $I_{\text{blur}}$ 为低通滤波结果，$\lambda$ 为锐化强度系数。

\paragraph{Mathematical Background}
算法上来讲，为了实现高效的滤波，在Separable Filtering中采用数学技巧降低计算量。本质上在某些核中可以使用线性技巧，将矩阵Kernel展开到$K = vh^{T}$的形式，然后只需要做两个一维乘法。

线性滤波的数学特征主要有以下两个：

\begin{itemize}
    \item Correlation: $$g(i,j) = \sum_{k,l}f(i+k,j+l)h(k,l)$$
    \item Convolution: $$g(i,j) = \sum_{k,l}f(i-k,j-l)h(k,l)$$
\end{itemize}

这两个算子都具备线性和平移不变性。

\begin{itemize}
    \item $h \circ (f_0+f_1) = h \circ f_0 + h \circ f_1$
    \item Shift Invarience: $(h \circ g)(i , j) = (h \circ f)(i + h , j + l) $ 
\end{itemize}

Filter和卷积神经网络中的特征提取算法有非常紧密的联系。

\subsubsection{Non-linear Filters}

\paragraph{Bilateral Filter}

Gaussian滤波的主要困难在于只考虑了像素位置的相近性，但是没有考虑像素灰度的差异，会忽略边缘，让重要的信息变模糊。

将Gaussian滤波（Domain Kernel）和颜色（Range Kernel）变化的结合起来，得到Bilateral Filter的权重

\[
w (i,j,k,l) = \exp(-\frac{(i-k)^2+(j-l)^2}{2\sigma_{d}^{2}}-\frac{||f(i,j) - f(k,l)||^2}{2\sigma_{r}^{2}})
\]

将此作为权重，结合softmax的归一化操作，得到滤波核

\[
g(i,j) = \frac{f(k,l)w(i,j,k,l)}{w(i,j,k,l)} \Rightarrow softmax(-\frac{(F_i-F_j)^2}{\sigma}f)
\]

这种滤波技术常用于美颜磨皮和ps技术, 多模态学习中的Kernel本质上就是可以学习的滤波器。

\paragraph{Joint Bilateral Filter}

如果有高清图片作为参考，进行图像滤波的时候可以直接参照正确答案查看颜色的相似性，实现去噪（用没Flash的参照给有Flash的进行去噪）。这也是神经网络中自注意力的灵感来源。如果自身的颜色细节是不好的，滤波器中的sigma参数应当依据无Flash的高清图片来确定，得到joint滤波。

上图是 No-Flash 图某一行的灰度剖面，$\sigma_d$ 控制参与平均的空间邻近范围；下图是 Flash 图同一行的剖面，$\sigma_r'$ 由 Flash 图确定颜色相似范围。两者分别对应权重公式中的域核（Domain Kernel）与值核（Range Kernel）。

技术路径分三步：

\begin{enumerate}[label=\arabic*.]
    \item 以 Flash 图为引导图（guide），对 No-Flash 图做 Joint Bilateral Filter，得到去噪后的基础层：
    \[
    I_{\mathrm{base}} = \mathrm{JBF}(I_{\mathrm{noflash}} \mid I_{\mathrm{flash}})
    \]
    \item 对 Flash 图做普通 Bilateral Filter 得到其平滑版本，再相除提取细节层：
    \[
    I_{\mathrm{detail}} = I_{\mathrm{flash}} / \mathrm{BF}(I_{\mathrm{flash}})
    \]
    \item 两者相乘，把 Flash 图的高频细节搬回低噪声的 No-Flash 图上：
    \[
    I_{\mathrm{result}} = I_{\mathrm{base}} \times I_{\mathrm{detail}}
    \]
\end{enumerate}

想要保持细节，则可以对原图进行滤波，提取高频细节。高频细节为了得到对比度，原图不带闪光灯饱和度低，对比时不宜使用减法提取，应该用除法提取对比度高的图像。

\subsection{UpSampling \& DownSampling}

\subsubsection{DownSampling and Anti-aliasing}
减小图像尺寸如果要通过降采样随机丢弃行列的方法进行实现，有可能出现走样的情况。这是因为采样频率无法反映图形中的高频噪声和高频细节，所以出现了不知所谓的结果。解决的方法是首先进行高斯滤波，将图像变得平滑。

将高斯平滑的方法和sampling排列起来，则可以构造出Gaussian金字塔，表现这个缩小的过程。

\begin{figure}[H]
    \centering
    \includegraphics[width=0.9\textwidth]{/home/simson/图片/截图/截图 2026-09-24 10-27-33.png}
    \caption{Gaussian 金字塔：每一级先做高斯模糊（Blur）再降采样（Subsample），逐级得到尺寸减半的图像}
\end{figure}

\paragraph{Laplacian Pyramid}

高斯金字塔每一级只保留低通信息，细节在模糊与降采样中被丢掉了。把相邻两级相减，就能把这些细节显式地取出来：

\[
L_i = G_i - \text{upsample}(G_{i+1}), \qquad G_i = L_i + \text{upsample}(G_{i+1})
\]

其中 $\text{upsample}(\cdot)$ 就是下一小节的插值放大（对应图中的 ``$\times 2$'' 节点），相减对应图中的圆圈；最粗的一级没有更粗的层可减，直接取 $L_4=G_4$。

\begin{figure}[ht]
    \centering
    \includegraphics[width=0.82\textwidth]{/home/simson/图片/截图/截图 2026-09-24 10-32-55.png}
    \caption{Laplacian 金字塔：$L_i$ 是 $G_i$ 减去放大后的 $G_{i+1}$ 所得的残差}
\end{figure}

右式说明这个表示是\textbf{可逆}的：从 $L_4$ 逐层加回去即可无损重建原图，而单用高斯金字塔则只有低通信息。$\{L_1,\dots,L_4\}$ 的优势可以顺着图看：

\begin{itemize}
    \item \textbf{分辨率分层}：$L_1$ 是最细的纹理与边缘，$L_2$、$L_3$ 依次是更粗的结构，$L_4$ 是那张最小的图，基本只剩低频。每一层对应一个明确的尺度，可以单独处理。
    \item \textbf{分频带（band-pass）}：每层只保留相邻两个尺度之间那一带频率。图中 $L_1$、$L_2$ 呈灰底加少量纹理，说明带通系数大多接近 0、非常稀疏，量化、阈值化和压缩都很划算（拉普拉斯金字塔是早期图像压缩的方案，也是小波变换的前身）。
    \item \textbf{冗余很小}：每层样本数按 $1/4$ 递减，加总约为原图的 $4/3$，只多花三分之一的空间就换来多尺度表示。
    \item \textbf{便于分尺度编辑}：亮度和细节被拆到了不同层，做色调映射、细节增强时可以只动某一层而不波及其他尺度；做图像融合时在不同层分别加权，还能避免拼接处的接缝。
\end{itemize}

图中每个 ``$\times 2$'' 节点都要靠插值把粗层放大来\textbf{预测}细层，预测得越准，残差层就越稀疏、越好压缩——这正是保边上采样（如 JBU）在实践中的价值。
Laplacian 金字塔因为逐级分频率保存信息，可以在数据传输中实现逐渐叠加，从模糊到清晰的效果。同时可以作为一种可以恢复原状的图像压缩技术。

下面总结并阐述了Laplacian金字塔的构造过程和处理流程：

\begin{figure}[H]
    \centering
    \includegraphics[width=0.70\textwidth]{/home/simson/图片/截图/截图 2026-09-24 10-41-19.png}
    \caption{Laplacian 金字塔的分解（左）与重建（右）流程}
\end{figure}

左半是分解：每级先用高斯滤波 $H_0$ 平滑、再降采样 $\downarrow 2$ 得到粗层，粗层经上采样 $\uparrow 2$ 与插值 $I$ 回到细层网格后，在求和点与原图相减（$-$）得到该级的细节残差，再进入下一级，如此逐级级联；右半是重建：各层数据经 $Q$ 后逐级上采样、插值并相加，最终得到输出 $O$。黄色块是各层的图像数据，蓝色块是可逆的滤波算子。

\paragraph{Mathematics of the Laplacian Pyramid}

根据构造流程，首先进行Gauss滤波，然后将不同程度的Gauss滤波进行做差和求微分过程：

\[
G(x, y ,\sigma) = \frac{1}{2 \pi \sigma} e ^{-\frac{x ^ 2 + y ^ 2}{2 \sigma^2}} \qquad \nabla^2 f = \frac{\partial^2 f}{\partial x ^ 2} + \frac{\partial^2 f}{\partial y^2}
\]

根据上述函数，将laplace算子作用在上面。

\[
\nabla^2 G(x , y ,\sigma) = (\frac{x ^ 2 + y ^ 2}{\sigma^4} - \frac{2}{\sigma^2})G(x, y, \sigma) = \frac{\partial G}{\partial \sigma} \frac{1}{\sigma}
\]

得到了Gauss函数非常重要的数学关系，使用展开式进行离散化。

\[
\sigma \nabla^2 G(x, y, \sigma) = \frac{\partial G}{\partial \sigma} \simeq (\frac{G(x,y,h\sigma) - G(x,y,\sigma)}{k\sigma - \sigma})
\]

化简得到：

\[
G(x, y ,k\sigma) - G(x, y, \sigma) \simeq (k-1)\sigma^2 \nabla^2 G
\]

这就是laplace金字塔的数学模型。

\subsubsection{UpSampling}

把小的图片放大，会出现非常模糊的情况，为了避免简单使用放大像素的方法，需要进行插值。

插值的数学原理可以用如下公式进行表示，以一元函数插值为例：

设采样点为整数位置 $x_k$，已知采样值 $f(x_k)$，插值即构造 $g(x)$ 满足 $g(x_k)=f(x_k)$，并用于估计采样点之间的函数值。记 $x$ 落在区间 $[x_k,x_{k+1}]$ 内，$t=x-x_k\in[0,1)$。

最近邻插值（Nearest Neighbor）：取离 $x$ 最近的采样值，$g(x)$ 是阶梯函数
\[
g(x)=f(x_k),\qquad k=\arg\min_{j}|x-x_j|
\]

线性插值（Linear）：用相邻两点连直线
\[
g(x)=(1-t)\,f(x_k)+t\,f(x_{k+1})
\]

Cubic 插值：在每个区间上用三次多项式，由两端函数值与一阶导数连续共四个条件定出系数
\[
g(x)\big|_{[x_k,x_{k+1}]}=a_3(x-x_k)^3+a_2(x-x_k)^2+a_1(x-x_k)+a_0
\]
\[
g(x_k)=f(x_k),\quad g(x_{k+1})=f(x_{k+1}),\quad g'(x_k^-)=g'(x_k^+),\quad g'(x_{k+1}^-)=g'(x_{k+1}^+)
\]

插值的方法有最近邻插值法，线性插值法。最近邻插值就是直接把最近邻居的数值搬过来，形成类似直方图的函数状态；比这个效果稍好的是线性插值，可以至少具有丰富的变化和导数。效果更好的是选取分段的高次函数，比如Cubic插值选用三次函数保证一阶连续。

放到图像上采样的插值算法中，双线性插值法是较为常用的，只需要知道周围四个邻居的数值，但是注意，双线性插值是一种非线性插值，结果中有交叉项。如果要追求好的效果，用Bicubic插值。

双线性插值（Bilinear）：用周围四个邻居。记 $u=x-i,\ v=y-j$，$(u,v)\in[0,1)^2$
\[
g(x,y)=(1-u)(1-v)f(i,j)+u(1-v)f(i+1,j)+(1-u)v\,f(i,j+1)+uv\,f(i+1,j+1)
\]
写成矩阵形式，交叉项 $uv$ 出现在展开后的二次项中：
\[
g(x,y)=
\begin{bmatrix}1-u & u\end{bmatrix}
\begin{bmatrix}f(i,j) & f(i,j+1)\\ f(i+1,j) & f(i+1,j+1)\end{bmatrix}
\begin{bmatrix}1-v\\ v\end{bmatrix}
\]

Bicubic 插值：用周围 $4\times4=16$ 个邻居，两个方向分别套用同一个三次核 $W$，可分离地写成
\[
g(x,y)=\sum_{m=-1}^{2}\sum_{n=-1}^{2} f(i+m,\ j+n)\,W(u-m)\,W(v-n)
\]
其中 $W$ 为三次卷积核（$a$ 通常取 $-0.5$）
\[
W(s)=
\begin{cases}
(a+2)|s|^3-(a+3)|s|^2+1, & |s|\le 1\\[2pt]
a|s|^3-5a|s|^2+8a|s|-4a, & 1<|s|<2\\[2pt]
0, & |s|\ge 2
\end{cases}
\]

\subsubsection{Joint Bilateral Upsampling}

\paragraph{Motivation: Why Upsample Jointly}

深度估计（立体匹配 stereo matching、多视图立体 multi-view stereo 等）需要为每个像素在一系列候选视差（disparity）上做匹配，代价随像素数增长，通常是整个流程中最贵的一步。如果最终只需要一张低分辨率的深度图，一个常见的加速方案是把流程拆成三步：

\begin{enumerate}[label=\arabic*.]
    \item 对输入图像降采样（downsample），得到低分辨率彩色图；
    \item 在低分辨率上估计深度图（estimate the depth image）；
    \item 把深度图上采样（upsample）回原始分辨率。
\end{enumerate}

前两步的代价大致与像素数成正比，因此把最贵的深度估计放到低分辨率上做，代价可以显著下降，这就是 ``for significant speedup'' 的含义。问题出在第三步：直接用上一节的最近邻、双线性或 Bicubic 插值放大深度图，并不能得到满意的结果。这些插值只关心样本之间的空间距离，完全不知道图像内容；而深度图在物体边界处存在跳变（discontinuity），插值核会把边界两侧的样本一起平均，于是边界被抹成一条斜坡，与真实的三维结构不符。最近邻插值虽然不会模糊，但会产生块状锯齿。

\begin{figure}[H]
    \centering
    \includegraphics[width=0.95\textwidth]{/home/simson/图片/截图/截图 2026-09-24 10-20-31.png}
    \caption{联合双边上采样的动机：降采样—估计深度—上采样，上采样时复用原始输入图像的边缘信息}
\end{figure}

于是自然的想法是：不要只在深度图内部插值，而是\textbf{复用原始输入图像中的边缘信息}（reuse the edge information in the original input image）。物体的边界在彩色图和深度图上出现在同一位置，高分辨率彩色图仍然保留着低分辨率深度图已经丢失的边界位置，因此可以把它作为引导图（guide），让插值权重在颜色发生跳变的地方被截断。

\paragraph{The JBU Formula}

Kopf 等人把上一节双边滤波的思想搬到上采样任务上，提出了联合双边上采样（Joint Bilateral Upsampling, JBU）。设低分辨率深度图为 $S$，高分辨率彩色引导图为 $\tilde I$，则高分辨率深度图上像素 $p$ 的取值为

\[
\tilde S_p = \frac{1}{k_p}\sum_{q_\downarrow\in\Omega} S_{q_\downarrow}\ f(\|p_\downarrow - q_\downarrow\|)\ g(\|\tilde I_p - \tilde I_q\|)
\]

其中各符号的含义为：

\begin{itemize}
    \item $p$：高分辨率输出图上的目标像素；$p_\downarrow$ 是它在低分辨率网格上的对应坐标（$\downarrow$ 表示低分辨率深度图中的坐标）。
    \item $q_\downarrow$：低分辨率深度图中落在 $p_\downarrow$ 邻域 $\Omega$ 内的样本点，求和遍历整个 $\Omega$。
    \item $S_{q_\downarrow}$：该样本点的低分辨率深度值，也就是被插值的信号。
    \item $\tilde I_p,\tilde I_q$：高分辨率彩色引导图在对应位置的取值。
    \item $f$ 与 $g$ 都是高斯函数（Gaussian function）：$f$ 作用在空间距离上，是域核（domain kernel）；$g$ 作用在颜色差上，是值核（range kernel）。
    \item $k_p$：归一化因子（normalization factor），即所有权重之和 $k_p=\sum_{q_\downarrow\in\Omega} f\,g$，使权重归一化到和为 1。
\end{itemize}

公式的结构与双边滤波完全一致，区别集中在两点：

\begin{itemize}
    \item 空间项 $f$ 是在\textbf{低分辨率}坐标上计算的，$\|p_\downarrow-q_\downarrow\|$ 度量的是样本到目标位置的距离。它扮演插值核的角色：距离越远权重越小，从而在稀疏的低分辨率样本之间平滑地重建出高分辨率的取值。
    \item 颜色项 $g$ 取的是\textbf{高分辨率}彩色图。当样本跨过物体边界时，$\|\tilde I_p-\tilde I_q\|$ 很大，$g$ 迅速趋近于 0，跨越边界的样本被排除在平均之外。
\end{itemize}

因此 $f$ 负责“插值”（保持平滑），$g$ 负责“保边”（在颜色跳变处截断），二者相乘得到一个形状随位置变化、由引导图决定的插值核。这正是它比固定核的 Bicubic 插值更贴合真实边界的原因。

\begin{figure}[H]
    \centering
    \includegraphics[width=0.95\textwidth]{/home/simson/图片/截图/截图 2026-09-24 10-22-53.png}
    \caption{联合双边上采样在一维信号上的过程：左列为信号（上：低分辨率深度；下：高分辨率彩色），中列为两个高斯权重 $f$、$g$，右列为相乘后的权重与上采样结果 $\tilde S_p$}
\end{figure}

上图把公式在一维信号上逐步展开。第一行是空间部分：低分辨率深度图在边界处有一个跳变，目标位置 $p$ 用蓝色星号标出，落在低的那个台阶上；空间权重 $f$ 是以 $p_\downarrow$ 为中心的高斯；它与信号相乘后，在低分辨率样本上得到一条以 $p$ 为中心的权重曲线。第二行是颜色部分：高分辨率彩色图在对应位置也有一个跳变，颜色权重 $g$ 在颜色与 $\tilde I_p$ 相同的区域接近 1，一跨过边界就掉到 0；于是把 $f$ 与 $g$ 相乘得到的组合权重中，位于边界另一侧的样本权重被压制为 0。最终结果是高分辨率深度图 $\tilde S_p$：既在物体内部平滑，又保持了锐利的边界，而且边界位置正好落在彩色图的边缘上。

\paragraph{Comparison with Other Upsampling Methods}

\begin{figure}[H]
    \centering
    \includegraphics[width=0.95\textwidth]{/home/simson/图片/截图/截图 2026-09-24 10-23-07.png}
    \caption{四种上采样方法在深度图边界处的对比（上排为放大的深度图区域，下排为对应的高分辨率彩色图）}
\end{figure}

同一张低分辨率深度图用四种方法放大，效果差别很明显：

\begin{itemize}
    \item \textbf{最近邻上采样（Nearest Neighbor）}：直接复制最近样本，深度值不产生新的中间值，但边界走样成块状的锯齿，呈阶梯状；
    \item \textbf{双线性 / Bicubic 上采样}：结果平滑，但深度边界被抹成一条逐渐过渡的斜坡，与彩色图中的锐利边缘不再对齐；
    \item \textbf{高斯上采样（Gaussian）}：本质上是各向同性的低通插值，模糊更严重，边界更加弥散；
    \item \textbf{联合双边上采样（JBU）}：在物体内部平滑过渡，在边界处与彩色图的边缘对齐，得到既连续又锐利的深度图。
\end{itemize}

\paragraph{Discussion}

\begin{itemize}
    \item 与上一节的联合双边滤波相比，这里的“联合”体现在两个不同分辨率的信号上：被滤波的是低分辨率深度图，提供颜色相似性判据的却是高分辨率彩色图。若引导图就是被滤波图像本身、且两者分辨率相同，公式即退化为普通的双边滤波。
    \item 计算上，每个高分辨率像素只需要遍历一个小邻域 $\Omega$（例如 $4\times 4$ 个低分辨率样本），比直接在高分辨率上重新估计深度便宜得多，因此适合作为“低分辨率估计 + 高分辨率恢复”这一套流程的最后一步。
    \item 这种用引导图决定滤波核形状的思路后来被推广为引导滤波（Guided Filter）等边缘保持滤波方法，也常出现在深度图增强、深度上采样以及需要从低分辨率标注恢复高分辨率预测的任务中。
\end{itemize}

Kopf, Johannes, et al. ``Joint bilateral upsampling.'' \textit{ACM Transactions on Graphics}, 2007.

\subsubsection{Image Blending}

这是Laplacian金字塔在图像混合中的重要应用。图像混合要让低频的噪声尽可能均匀混合，同时体现出高频部分的线条，物体等等，laplace金字塔可以分开不同的频率，因此常常被用在类似的操作中。

具体做法分四步：

\begin{enumerate}[label=\arabic*.]
    \item 对两张源图分别建立 Laplacian 金字塔 $\{L_i^{A}\}$、$\{L_i^{B}\}$；
    \item 对掩膜 $M$ 建立 Gaussian 金字塔 $\{M_i\}$，逐级模糊相当于逐级加宽过渡带；
    \item 在每一层分别混合 $L_i = M_i \odot L_i^{A} + (1-M_i)\odot L_i^{B}$；
    \item 把混合后的各层折叠（collapse，逐级上采样并相加）回一张图像。
\end{enumerate}

\begin{figure}[H]
    \centering
    \includegraphics[width=0.72\textwidth]{/home/simson/图片/截图/截图 2026-09-24 11-04-41.png}
    \caption{Laplacian 金字塔融合：每一行是一个频带，两张图在该频带上按掩膜加权相加，最后把所有频带折叠回图像}
\end{figure}

图中每一行就是一个频带，自上而下由细到粗：最上面一行是最细的高频带，掩膜几乎还是硬边；越往下频带越粗，掩膜被模糊得越厉害，最下面一行只剩大范围的平滑过渡，混合的主要是颜色与亮度。这正是图中 Key Idea 的含义——``blend lower frequency bands over larger spatial ranges, high frequency bands over small spatial ranges''：低频用大范围过渡，所以看不到接缝；高频用小范围过渡，所以纹理和边缘不会被糊掉。直接按掩膜逐像素加权会在边界留下可见的接缝，而先把掩膜模糊又会连带糊掉细节，金字塔融合用“过渡带宽度随频带变化”同时避开了这两个问题。

这也点出了金字塔表示的价值：同一个边界在不同频带上需要不同的空间尺度，只有先把图像拆成频带，才可能对它们分别处理（Burt \& Adelson, 1983）。

\subsection{Image Transformation}

\paragraph{Image Warping}

Image Warping指的是对图像的几何变形，比如旋转，伸缩，拉伸，变形等等。要点在于如果直接映射原来图像到新图像，则会出现空洞或重叠现象。重点方法在于扫描目标图像，逆映射到原图的实数值，然后通过插值方法取得像素的数值，保证不会重复计算，全面覆盖。

\begin{figure}[H]
    \centering
    \includegraphics[width=0.75\textwidth]{/home/simson/图片/截图/截图 2026-09-24 11-23-04.png}
    \caption{反向映射：对目标图的每个像素 $x'$ 用逆映射求出源位置 $x=\hat h(x')$，再在源图上重采样}
\end{figure}

对目标图像 $g(x')$ 中的每一个像素：

\begin{enumerate}[label=\arabic*.]
    \item 用逆映射 $\hat h$ 算出它在源图中的位置 $x=\hat h(x')$；
    \item 在源图 $f$ 的该位置重采样（resample），把结果复制到 $g(x')$。
\end{enumerate}

两个实现细节：

\begin{itemize}
    \item $\hat h$ 是 $h$ 的逆映射。如果反过来遍历源像素、算出目标位置再写过去，目标图上会出现空洞（相邻源像素被拉开）与重叠（多个源像素落到同一格）；反向映射则保证每个目标像素恰好被写入一次。
    \item 重采样用双线性插值：$\hat h(x')$ 通常不是整数坐标（图中下排的红点落在四个格点之间），需要由它周围四个源像素加权得到取值，也就是 3.3.2 节讲的双线性插值。
\end{itemize}

\paragraph{Image Morphing}

Image Morphing指的是两个图像进行融合叠加的方法，实质上是分块的Warping和blending。

Morphing 的具体步骤（源图 $s$、目标图 $u$，进度参数 $\mu\in[0,1]$）：

\begin{enumerate}[label=\arabic*.]
    \item 在两幅图上选取一一对应的特征点（feature points），各自做 Delaunay 三角剖分，得到网格 $T_s$、$T_u$；
    \item 对三角形的位置做线性插值，得到中间网格 $T = \mu T_s + (1-\mu) T_u$；
    \item 按 $T_s \to T$ 与 $T_u \to T$ 的对应关系，在每个三角形内把两幅图分别 warp 成 $\tilde s$、$\tilde u$（三角形内部的映射由三个顶点唯一确定，是仿射变换）；
    \item 把两幅已经对齐的图像按颜色混合，输出 $o = \mu\tilde s + (1-\mu)\tilde u$。
\end{enumerate}

\begin{figure}[H]
    \centering
    \includegraphics[width=0.75\textwidth]{/home/simson/图片/截图/截图 2026-09-24 11-23-21.png}
    \caption{Image morphing：$T_s$、$T_u$ 是源图与目标图的 Delaunay 三角剖分，插值出中间网格后在每个三角形内 warp 并混合}
\end{figure}

$\mu$ 由 0 变到 1 时，输出的几何形状与颜色同时从源图过渡到目标图：前者靠 warping，后者靠 cross-dissolve（线性混合），morphing 正是两者的组合。

